\section{Quản lí nhân viên sử dụng hệ thống}

Nhóm tính năng ``Quản lí nhân viên sử dụng hệ thống" đóng vai trò quan trọng trong việc đảm bảo hoạt động của hệ thống được vận hành hiệu quả và an toàn trong bệnh viện. Thông qua nhóm tính năng này, quản trị viên có thể kiểm soát được vùng giới hạn thông tin mà nhân viên bệnh viện sẽ được truy cập, từ đó đảm bảo mỗi cá nhân có quyền truy cập đúng để phù hợp với vai trò và trách nhiệm của mình.

Nhóm tính năng này đặc biệt quan trọng trong môi trường bệnh viện, nơi có sự tham gia của nhiều bộ phận như bác sĩ, y tá, điều dưỡng, nhân viên tiếp nhận, và cán bộ quản lí. Việc quản lí chặt chẽ nhân viên sẽ giúp ngăn ngừa truy cập trái phép, giảm thiểu sai sót trong xử lý bệnh án và tăng cường bảo mật thông tin.

Nhóm tính năng bao gồm các tính năng cụ thể như:

\begin{itemize}
  \item Thêm nhân viên mới
  \item Cập nhật thông tin nhân viên
  \item Phân quyền nhân viên
  \item Tra cứu thông tin nhân viên
\end{itemize}

\subsection{Thêm nhân viên mới}

\noindent \textbf{Giới thiệu tính năng}

Tính năng này cho phép quản trị viên hệ thống thêm nhân viên mới (là nhân viên bệnh viện bao gồm nhân viên y tế, nhân viên hành chính, ...) vào hệ thống và phân quyền để họ có thể truy cập và sử dụng các chức năng phù hợp với vai trò của mình.

\noindent \textbf{Các yêu cầu chức năng}

\begin{itemize}
  \item \textbf{REQ\_01} Hệ thống phải tự động sinh mã nhân viên duy nhất cho mỗi nhân viên mới được thêm vào.
  \item \textbf{REQ\_02} Hệ thống phải kiểm tra tính hợp lệ của các trường bắt buộc bao gồm họ tên, ngày sinh, số CCCD, giới tính, ... trước khi tạo tài khoản nhân viên mới.
  \item \textbf{REQ\_03} Hệ thống phải tạo được tài khoản nhân viên mới và lưu trữ thông tin cá nhân đã được cung cấp vào cơ sở dữ liệu.
  \item \textbf{REQ\_04} Hệ thống phải tự động sinh tên đăng nhập và mật khẩu mặc định.
\end{itemize}

\subsection{Phân quyền nhân viên}

\noindent \textbf{Giới thiệu tính năng}

Tính năng này cho phép quản trị viên hệ thống gán quyền truy cập hệ thống cho nhân viên bệnh viện dựa trên vai trò của họ, nhằm đảm bảo mỗi người chỉ có thể sử dụng các chức năng phù hợp với phạm vi trách nhiệm và nhiệm vụ của mình.

\noindent \textbf{Các yêu cầu chức năng}

\begin{itemize}
  \item \textbf{REQ\_01} Hệ thống phải hỗ trợ các vai trò nhân viên mặc định với các quyền truy cập được cấu hình sẵn.
  \item \textbf{REQ\_02} Hệ thống phải từ chối các truy cập mà nhân viên không có quyền truy cập.
\end{itemize}

\subsection{Cập nhật thông tin tài khoản nhân viên}

\noindent \textbf{Giới thiệu tính năng}

Tính năng này cho phép quản trị viên hệ thống chỉnh sửa, cập nhật thông tin cá nhân và thông tin hệ thống của nhân viên bệnh viện đã được thêm vào hệ thống trước đó, đảm bảo dữ liệu luôn chính xác, đầy đủ và kịp thời.

\noindent \textbf{Các yêu cầu chức năng}

\begin{itemize}
  \item \textbf{REQ\_01} Hệ thống phải cho phép chỉnh sửa các trường thông tin cá nhân như họ tên, ngày sinh, giới tính, số CCCD, số điện thoại, email...
  \item \textbf{REQ\_02} Hệ thống phải kiểm tra tính hợp lệ của thông tin được cập nhật trước khi lưu.
\end{itemize}

\subsection{Tra cứu thông tin nhân viên}

\noindent \textbf{Giới thiệu tính năng}

Tính năng này cho phép quản trị viên tìm kiếm và xem thông tin chi tiết của nhân viên trong hệ thống, phục vụ cho việc quản lý, hỗ trợ hoặc xác minh thông tin khi cần thiết.

\noindent \textbf{Các yêu cầu chức năng}

\begin{itemize}
  \item \textbf{REQ\_01} Hệ thống phải hiển thị danh sách nhân viên một cách rõ ràng, có phân trang nếu kết quả quá nhiều.
  \item \textbf{REQ\_02} Hệ thống phải cho phép xem thông tin chi tiết của từng nhân viên khi chọn vào một nhân viên cụ thể.
\end{itemize}
